\documentclass[a4paper, 12pt]{book} % changed from scrbook
% not using \documentclass[a4paper, twoside=semi]{scrreprt} anymore

\usepackage{InfoSec}

\definecolor{bg}{RGB}{20,20,20}
\definecolor{fg}{RGB}{220,220,220}
\definecolor{keyword}{RGB}{130,170,255}
\definecolor{comment}{RGB}{120,120,120}
\definecolor{string}{RGB}{255,200,120}

\lstdefinestyle{darkbash}{
    backgroundcolor=\color{bg},
    basicstyle=\scriptsize\ttfamily\color{fg},
    keywordstyle=\color{keyword},
    commentstyle=\color{comment},
    autogobble=true,
    stringstyle=\color{string},
    breaklines=true,
    showstringspaces=false,
    frame=single
}

% \usepackage{titlesec}
% \setkomafont{chapter}{\Huge\bfseries}
% \setkomafont{section}{\Large\bfseries}
% \setkomafont{subsection}{\large\bfseries}
\usepackage{fancyhdr, lastpage}
\usepackage{lstautogobble}
\usepackage{euscript}


\makeatletter
\providecommand{\subtitle}[1]{\gdef\@subtitle{#1}}
\providecommand{\@subtitle}{}
\makeatother


\usepackage{booktabs}
\usepackage{array}
\usepackage{multirow}




\usepackage{newfloat}
%\usepackage{tocloft}

% % Define new float environment
% \DeclareFloatingEnvironment[
%   fileext=lof,
%   listname={List of Flowcharts},
%   name=Flowchart,
% ]{flowchart}


\pagestyle{fancy}
%\fancyhf{}
\lhead{}
\chead{}
%\rhead{}
\cfoot{\thepage\ of \pageref{LastPage}}
\rfoot{Version 8.3}




\RequirePackage{etex} 


%--- The documentation for the title can be found at: https://en.wikibooks.org/wiki/LaTeX/Title_Creation
\title{Information Security \\ Public Key Infrastructure}
\author{Ahmed U. Ashraf}
\date{\today }
% \subtitle{Version~\version}

\begin{document}

% \include{glossary}


\maketitle
%--- we would like to hyperlink each item in contensts page
% \let\hypercontentsline=\contentsline
% \renewcommand{\contentsline}[4]{\hypertarget{toc.#4}{}\hypercontentsline{#1}{#2}{#3}{#4}}
%--- We print ToC: Table of Contents

\DeclareFontShape{U}{stmry}{b}{n}{<-> ssub * stmry/m/n}{}

\tableofcontents\label{toc}
%--- We print the list of figures, this might go at the end.




% \listoffigures
% %--- We print the list of tables, this might also go at the end.
% \listoftables
% %--- This is to add list of symbols as glossary
% % \printunsrtglossary[type=symbols,style=long,title={List of Symbols}]


% % List of Algorithms (- taken the idea from Discrete Algorithmic Mathematics by Stephen B.Maurer and ANthony Ralston.pdf)

% \listofalgorithms
% \addcontentsline{toc}{chapter}{List of Algorithms}

% \listofflowcharts
% %\addcontentsline{toc}{chapter}{List of Flowcharts}







\chapter{Introduction}

Public-key cryptography is a cryptographic technique that enables entities to securely communicate on an insecure public network, and reliably verify the identity of an entity via digital signatures. A \emph{public key infrastructure}\footnote{This includes Hardware, software and policies.} (PKI) is a system for the creation, storage, management, distribution and revokation of public keys and digital certificates, which are used to verify that a particular public key belongs to a certain entity. The PKI creates digital certificates that map public keys to entities, securely stores these certificates in a central repository and revokes them if needed. PKI consists of
\begin{enumerate}
    \item A \emph{certificate authority} (CA) which stores, issues, and signs digital certificates.
    \item A \emph{registration authority} (RA) which verifies the identity of entities requesting their digital certificates to be stored at the CA.
    \item A \emph{central directory}, a secure location in which keys are stored and indexed.
    \item A \emph{certificate management system}, which manages things like the access to stored certificates or the delivery of the certificates to be issued;
    \item A \emph{certificate policy}, which states the PKI's requirements concerning its procedures. Its purpose is to allow outsiders to analyze the PKI's trustworthiness.
\end{enumerate}


\section{Digital certificates}

Certificates are designed to bind a subject to the subject's public key. These certificates should be digitally signed by another entity called \emph{Certificate Authority} (CA).

\begin{example}
    OpenSSL, Let's Encrypt are examples of CA.
\end{example}


In order to get the certificate from the server \emph{google.com} of the host \emph{google.com}.
\begin{lstlisting}[style=darkbash, language=bash]
echo | \
openssl s_client \
-connect google.com:443 \
-servername google.com | \
openssl x509 \
-outform PEM \
> google_cert.pem
\end{lstlisting}


What does X.509 certificate contains:
\begin{table}[H]
    \begin{tabular}{|c|c|} \hline \\
        Version & $\texttt{3 (0x2)}$\\ \hline 
        Serial Number & $\texttt{\tiny 13:9c:df:29:a8:b5:bc:52:12:41:3e:ee:af:d1:8c:be}$ \\ \hline 
        Signature Algorithm & $\texttt{sha256WithRSAEncryption}$ \\ \hline 
        Issuer &  C=US, O=Google Trust Services, CN=WR2 \\
        Validity & \\
        Not Before & May 25 08:36:19 2026 \\
        Not After & Aug 17 08:36:18 2026 \\
        Subject & CN=*.google.com \\
        Subject Public Key Info &  \\
        Public Key Algorithm &  id-ecPublicKey \\
        NIST CURVE &  P-256 \\
        Public-Key & (256 bit) \\
    \end{tabular}
\end{table}

\section{Certificate Signing Requests (CSRs)}












% \chapter{References}

%-- Computation Engineering by Ganesh Gopalakrishnan

\cite{adamson2012set}

% \nocite{*}



\bibliography{references}
\bibliographystyle{amsalpha}

\printindex




\end{document}